\documentclass[11pt, letterpaper]{article}
\usepackage[letterpaper,margin=1in]{geometry}
\usepackage{amsmath}
\usepackage{amssymb}
\usepackage{bigints}
\linespread{1.25}

\newcommand{\X}{\mathcal{X}}
\newcommand{\Y}{\mathcal{Y}}
\renewcommand{\d}{\textrm{d}}

\begin{document}
An Automated Market Maker (AMM) is a protocol which permits the execution of buy and sell orders on a blockchain.

Uniswap v3 is a constant product market maker, meaning that it maintains invariant function $\mathcal{I}(x,y)=x\cdot y =L^2$ where $L$ denotes the liquidity $x$ and $y$ are virtual reserves of tokens $\X$ and $\Y$, respectively. A virtual reserve for token $\X$ is defined as $x=x_{\rm offset} + x_r$ where $x_r$ is the real reserves, and $x_{\rm offset}$ are some other quantities.
A liquidity provider (LP) is an actor who transfers tokens into liquidity pools and are rewarded by the fees paid by traders who use that pool to swap tokens. 
In Uniswap v3, LPs wishing to provide liquidity must specify a price range over which their tokens will be used for trading activities.
Here, a price range defined as $R:=[p_l,p_u)$ where $p_l=\frac{x_l}{y_l}$ and $p_u=\frac{x_u}{y_u}$ must be specified in ticks $\tau$. Specifically, price $p(\tau)=\Big(\beta^{(p)}\Big)^\tau$ where $\beta^{(p)}=1.0001$ and $\tau$ is a multiple of $\delta_\pi$ (a value determined by swap fees $\phi=0.05\%,0.3\%,1\%$). 
The authors define a unitary price range as the range spanned by two consecutive ticks $[p_l, p_u)$ such that $p_l(\delta_\pi\cdot i)$ and $p_u(\delta_\pi\cdot (i+1))$.

Consider a square-root price range $[\pi_l,\pi_u):\pi_l=\sqrt{p_l},\, \pi_u=\sqrt{p_u}$ that is in the form $R_1\cup\ldots\cup R_n: n\geq 1$.
For $i\leq n,\, R_i=[\pi_{l_i},\pi_{u_i})$ contains an amount $L_i$ of liquidity, and for $i<n,\,\pi_{u_i}=\pi_{l_{i+1}}$ i.e. the ranges are consecutive.
An LP who deposited (resp. withdrew) the same amount of liquidity $\Delta L$ on range $R_i$, where $L_i+\Delta L\geq 0\forall i$, will deposit (resp. withdraw) the following amount of tokens into the pool:
\begin{equation}
\label{eq:liq}
\Delta x_r = \Delta L \cdot (\frac{1}{\pi_0^R}-\frac{1}{\pi_u})\quad \textrm{and}\quad \Delta y_r = \Delta L\cdot (\pi_0^R-\pi_l)
\quad \textrm{s.t.}\quad
\pi^R = \begin{cases}
    \pi_u & \text{if } \pi_u < \pi \\
    \pi_l & \text{if } \pi < \pi_l \\
    \pi   & \text{otherwise}
\end{cases}
\end{equation}
In this case, the liquidity provider is said to have deposited an amount of liquidity $\Delta L$ onto price range $R$.

If an LP wishes to withdraw (burn) liquidity $\Delta L$ on price range $R$, then he would retrieve $\Delta x_r$ and $\Delta y_r$ real reserve tokens as specified in \eqref{eq:liq}, and in addition earn $\Delta x_{\rm fee}=\Delta L \cdot (\Phi_R^\X(t)-\Phi_R^\X(t_0))$ and $\Delta y_{\rm fee}=\Delta L \cdot (\Phi_R^\Y(t)-\Phi_R^\Y(t_0))$
Here, accumulators $\Phi_R^\X$ and $\Phi_R^\Y$ are updated at every transaction. Suppose $a_i$ tokens $\X$ are transfered into the pool in exchanged for tokens $\Y$ via a swap operation, then $\Phi_R^\X\leftarrow\Phi_R^\X + \frac{a_i}{1-\phi}\cdot\frac{\phi}{L}$, where $L$ is the available liquidity on range $R$.
$\frac{a_i}{1-\phi}$ represents the input tokens actually used in the swap, which implies that $\frac{a_i\phi}{1-\phi}$ represents the accumulated fees from that operation and that $\frac{a_i}{1-\phi}\cdot\frac{\phi}{L}$ represents the amount of accumulated fees per unit of liquidity on the range $R$.

\subsection*{LP Strategy}
However, if a liquidity provider withdraws their tokens at a time where the pool price is significantly different from the one at deposit time, the value of their tokens net of fees will be lower than the value of their original tokens if they had not been deposited into the pool. This loss is only materialized when LPs withdraw their tokens from the pool, which is called \textit{Impermanent Loss} (IL) \cite{doi:10.1137/23M1606149}. 

Consider two strategies: (i) a liquidity providing stategy where tokens are deposited into a pool, and (ii) a strategy where tokens are withheld from the pool under $\bf{H}_0:$ Tokens $\X$ and $\Y$ outside the pool can be valued using pool price.. The value of strategies (i) and (ii) in reference num\'eraire token $\Y$ is denoted $V_P$ and $V_H$, respectively.Suppose that an amount of liquidity $\Delta L$, corresponding to token quantities $\Delta x_r$ and $\Delta y_r$ via equation \eqref{eq:liq}, is added to \textit{unitary} square-root price range $R:=[\pi_l,\pi_u)$ at time $t_0$. Let $p_i=\pi_i^2$ denote the price at time $t_i$, then at later time $t_1$,
\begin{equation}
V_P(t_1)=\Delta L\cdot \Bigl(\Bigl(\frac{1}{\pi_1^R}-\frac{1}{\pi_u}\Bigr)\cdot\pi_1^2+\bigl(\pi_1^R-\pi_l\bigr)\Bigr)
\end{equation}
Now suppose that quantities $\Delta x_r$ and $\Delta y_r$ are held outside of the pool (and under $\bf{H}_0$ are priced at $p_0=\pi_0^2$ at time $t_0$). 
\begin{equation}
V_H(t_1)=\Delta L\cdot \Bigl(\Bigl(\frac{1}{\pi_0^R}-\frac{1}{\pi_u}\Bigr)\cdot\pi_1^2+\bigl(\pi_0^R-\pi_l\bigr)\Bigr)
\end{equation}

These results above can be generalized to a full liquidity curve $(\Delta L_i)_i$ spread over unitary range $[\pi_{l_i},\pi_{u_i})$. 
For a given square-root price $\pi$, $\Delta L_\pi = \Delta L_i$ for $\pi\in[\pi_{l_i},\pi_{u_i})$, then the function $\pi\mapsto\Delta L_\pi$ is piecewise constant and right-continuous with a left limit, allowing for integrability w.r.t $\pi$

Consider a liquidity provider adding a liquidity curve $(\Delta_\pi)_\pi$ to the pool while keeping $x_0$ tokens $\X$ and $y_0$ tokens $\Y$ outside the pool (under $\bf{H}_0$). Assume the integrability condition 
\begin{equation}
\bigintsss_1^\infty \frac{\Delta L_\pi}{\pi^2}d\pi + \bigintsss_0^1 \Delta L_\pi d\pi < \infty
\end{equation}
At time $t_0$, when the price in the pool is $p_0$, their tokens admit the following $\Y$-value:
\begin{equation}
V_P(t_0)= x_0\cdot p_0 + y_0 + p_0\cdot\bigintsss_1^{+\infty} \frac{\Delta L_\pi}{\pi^2}d\pi + \bigintsss_0^{\pi_0} \Delta L_\pi d\pi < \infty
\end{equation}
At time $t_1$, when the price in the pool is $p_1$, their tokens $\Y$-value net of swap fees is:
\begin{equation}
\label{eq:strat_payoff}
V_P(t_1)= x_0\cdot p_1 + y_0 + p_1\cdot\bigintsss_1^{+\infty} \frac{\Delta L_\pi}{\pi^2}d\pi + \bigintsss_0^{\pi_0} \Delta L_\pi d\pi < \infty
\end{equation}
So, an LP depositing a liquidity curve implies a payoff (net of fees) which is concave as a function of price (see equation \eqref{eq:strat_payoff}).
From here the authors calculate the so-called ``Greeks'': (Delta) $\frac{\partial V_P(t_0)}{\partial p_0}$ and (Gamma) $\frac{\partial V_P(t_0^{\pm})}{\partial p_0}$ which gives the sensitivity of the strategy's global position $V_P(t_0)$ with respect to the $X-Y$ exchange rate. 
Next, the authors make use of the Greeks and the Carr-Madan formula\cite{ctx60986673600002466} to show that any smooth concave payoff function can be generated by an LP strategy.

\subsection*{Fees Modeling}
Let latent price process $(p_t)_t$ be defined on a filtered probability space $(\Omega, \mathcal{F}, (\mathcal{F})_{0\leq t\leq T},\mathbb{P})$.
Let hypothesis $\bf{H}_p:$ the latent price process follows an It\^o dynamics of the form: $\d p_t / p_t= \mu_t\d t + \sigma_t\d W_t$ with progressively measurable drift $(\mu_t)_t$ and volatility $(\sigma_t)_t$: both $\mu$ and $\sigma$ are bounded.

Assuming that token swaps mainly occur because of arbitrageurs between the pool and another reference exchange, swap trades will occur when the latent price process $(p_t)_t$ changes by $\gamma\%$ where $\gamma=\phi$ is the swap fee paid out to liquidity providers.
In other words, the swap times $(t_i)_i$ are described by $t_{i+1}=\inf\{t\ge t_i:p_t=p_{t_i}(1+\gamma)^{\pm 1}\}$.

Then, the amount of fees in $\X$ and $\Y$ accumulated over the period $[0,T]$ is defined as:
\begin{equation}
\label{eq:fees}
\text{Fees}^{\X}_{0\rightarrow T} = \sum_{R=[\pi_l,\pi_u)} \Delta L_\pi\cdot \bigl(\Phi_R^\X(T)-\Phi_R^\X(0)\bigr)
\quad \text{and}\quad
\text{Fees}^\Y_{0\rightarrow T} = \sum_{R=[\pi_l,\pi_u)} \Delta L_\pi\cdot \bigl(\Phi_R^\Y(T)-\Phi_R^\Y(0)\bigr)
\end{equation}

\newpage
\bibliographystyle{plain}
\bibliography{IPR_}

\end{document}
